% !PRINTED 332
\weberpage{332}
%<<CONT
2. \gesperrt{If we call the form $T$ the form belonging to the basis $(\alpha_\nu)$ of $\mathfrak{a}$, then the forms belonging to different positive bases of the same ideal are equivalent.}

Let us denote the conjugate values of $\alpha_s, \omega_s$ by

\[
\alpha_{s,1},\ \alpha_{s,2},\ \dots,\ \alpha_{s,n};\quad \omega_{s,1},\ \omega_{s,2},\ \dots,\ \omega_{s,n}
\]

and set

\weq{10}{\alpha_{s,r} = \overset{\nu}{\Sigma}\, a_{s,\nu}\,\omega_{\nu,r},}

then there results

\weq{11}{(\Sigma \pm \alpha_{1,1}\,\alpha_{2,2}\,\dots\,\alpha_{n,n})^2 = N(\mathfrak{a})^2 \varDelta,}

where $\varDelta$ is the fundamental number of the field $\Omega$, namely the determinant square:

\[
\varDelta = (\Sigma \pm \omega_{1,1}\,\omega_{2,2}\,\dots\,\omega_{n,n})^2
\]

(Bd.\,II, \S\,162).

If $\lambda_1, \lambda_2, \dots, \lambda_n$ are the conjugate values of $\lambda$, then by (4)

\weq{12}{\lambda_r = \overset{\nu}{\Sigma}\, \alpha_{\nu,r}\, t_\nu}

is a linear substitution for the variables $t_\nu$ with the substitution determinant

\weq{13}{r = N(\mathfrak{a})\,\sqrt{\varDelta},}

and through this the form $T$ passes, by (5), into the product

\weq{14}{T' = [N(\mathfrak{a})]^{-1}\,\lambda_1\,\lambda_2\,\dots\,\lambda_n}

since $N(\lambda)$ is the product of the conjugate values of $\lambda$ (Bd.\,II, \S\,151). We can apply invariant theory (Bd.\,I, \S\,65) to this, and if we form the \gesperrt{Hesse} determinant:

\weq{15}{H = \Sigma \pm \frac{\partial^2 T}{\partial t_1 \partial t_1}\,\frac{\partial^2 T}{\partial t_2 \partial t_2}\cdots\frac{\partial^2 T}{\partial t_n \partial t_n},}

then

\weq{16}{H' = r^{-2}\,H.}

If, however, we form $H'$ by (14) in the variables $\lambda_\nu$, and observe that an $n$-rowed determinant in which the diagonal terms are $=0$ and all the other terms are $=1$ has the value $(-1)^{n-1}(n-1)$\footnote{One can easily prove this theorem by extending the determinant through the addition of an $(n+1)$th row and an $(n+1)$th column, in which the added row contains only \gesperrt{ones}, and the column contains zeros with the exception of the last element.}, then it follows:
%CONT>>
